\subsection{物理约束分析}
\label{sec:discussion-physics}

物理约束通过将 Lamb 波传播的支配物理规律直接嵌入训练目标来嵌入领域知识。与将信号视为像素级回归问题的通用均方误差(MSE)损失不同，$\mathcal{L}_{\text{physics}}$ 中的三个物理项（第~\ref{sec:physics-loss} 节）对诊断特异性声学特征的偏差进行惩罚：到达时间、峰值振幅和波形形态。到达时间保留对首波检测中的误差进行惩罚，直接影响通过飞行时间测量的缺陷定位精度。振幅保留对峰值信号振幅的畸变进行惩罚，这对于缺陷尺寸评估至关重要，因为缺陷界面的声能反射率随缺陷严重程度缩放。波形形态通过 Lin 一致性相关系数(CCC)应用于信号包络来 enforcement，惩罚可能掩盖与复杂缺陷几何相关的模态转换效应的形状畸变。

均方误差(MSE)对于此应用失效，因为它优化聚合波形相似性，隐含地倾向于过度平滑的输出，抑制包括 arrival-time discontinuities 在内的高频瞬态。第~\ref{sec:results-ablation} 节的消融结果证实了这一机制：移除 $\mathcal{L}_{\text{physics}}$ 同时保持具有竞争力的信噪比(SNR)（$12.1$ vs $12.4$ dB），会导致到达时间误差(TOA)从 $0.42$ 增加到 $0.68$ $\mu$s，振幅误差从 $3.1\%$ 上升到 $7.4\%$，F1 分数从 $0.91$ 降至 $0.79$。因此，物理约束作为归纳偏置发挥作用，引导网络朝向物理上合理的解，同时不损害去噪能力。

一个关键的实际后果是泛化优势。因为物理项编码的是通用声波物理而非特定试样的信号特征，所以学习到的特征保留行为可以在不同试样几何形状间转移。这解释了第~\ref{sec:results-generalization} 节中报告的 intact-to-defect 和跨厚度泛化结果，其中使用物理损失的 PMDF-Net 在缺陷和厚度偏移测试集上保持了近乎 intact 的性能。